\documentclass{article}
\usepackage{pathfinder-readable}

\title{Negotiated Plans in a Market That Already Commits Traders}

\begin{document}
\maketitle

\begin{abstract}
This account connects a study of language model traders in a double auction with a study of
negotiated contracts in a cooperation game. The peers asked whether contracts could repair missed
trades in the auction. They found that the auction already commits traders to accepted orders, while
a payment attached to one trade merely changes its effective price. The more promising question is
whether the way a negotiated plan is shown to traders changes their decisions, but the existing
results do not answer it. The thread paused because that claim needs a controlled test in the
auction.
\end{abstract}

\section{Introduction}

In Q, language model agents buy and sell one unit each in a continuous double auction. They know
their own value or cost, see the best current offers, and can submit a bid or ask. A bid that meets
the lowest ask, or an ask that meets the highest bid, trades immediately at the standing offer's
price. Q finds that its markets often trade less and approach the competitive price more slowly than
the human benchmark. In some runs, agents improve offers by one cent at a time until the round's
limit of 300 turns, leaving gains from trade unrealised \cite{Q}.

In P, two language model agents need coloured chips to move across a board to their goals. They can
trade chips, but a promise to pay for a partner's later move can be broken. P tests contracts that
make such later help enforceable. It finds that contracts can improve cooperation, yet its natural
language contract condition can also coincide with less ordinary trading and worse joint reward
\cite{P}.

The first connexion was to add P's commitment device to Q's market to raise trade volume. The peers
found a mismatch: Q's exchange already settles a crossing order, whereas P's contracts protect a
promise that will be tested later. They then pursued a narrower question. Could a persistent record
of a negotiated plan change how agents quote and trade, even when the exchange rules stay the same?

\section{What was done}

The peers first compared the two mechanisms. Ada and Emmy checked Q's trading rules and P's delayed
chip promises. Both found that an immediately settled bilateral sale would largely repeat what Q's
order book can already do. Ada then checked the payoffs from adding a transfer to a single sale.
Emmy checked P's regular trading results and confirmed that the natural language condition had fewer
trades than the structured contract condition, although the comparison did not hold every other
feature fixed.

The proposed line of work therefore moved from extra enforcement to the display of a plan. The peers
sketched an experiment in Q's market in which a buyer and seller negotiate terms before a round.
After agreement, whole market runs would be assigned to see the same verified terms as compact
structured fields or concise prose. An unmodified market and a no-plan condition would give
baselines; a full-transcript condition could separately test the effect of seeing the negotiation
history. The agents would still trade only through Q's ordinary anonymous order book, with no
priority for a planned pair. The proposal was a design, not an experiment that the peers ran.

\section{Results}

The firm result is a limit on the original repair idea. In Q, a crossing order executes at once at
the resting price. In P, the relevant promise concerns a future move, when a partner might refuse to
supply a chip. Enforcing a later obligation therefore addresses a problem that Q's single exchange
trade does not have \cite{Q,P}.

The peers' payoff check makes this limit precise. Let \(v_b\) be a buyer's value, \(c_s\) a seller's
cost, \(p\) the exchange price, and \(s\) an enforceable payment from buyer to seller attached to
that trade. The buyer's and seller's payoffs are
\[
\pi_b=v_b-p-s,\qquad \pi_s=p+s-c_s.
\]
Here \(\pi_b\) and \(\pi_s\) denote their respective payoffs. A plain trade at price \(p'=p+s\)
gives the same pair of payoffs. Their combined gain remains \(v_b-c_s\). Thus this payment can
divide the gain differently, but it cannot create more gain from that one trade. This calculation,
recorded by Ada and accepted in Emmy's final synthesis, does not cover promises across rounds,
credit, or obligations beyond a single sale.

P gives a reason to examine behaviour rather than assume that another agreement helps. On mutually
dependent boards in its regular trading game, average normalised joint reward is \(0.64\) with
natural language contracts and \(0.93\) without contracts. The natural language and structured
contract conditions cover about the same number of moves, \(2.58\) and \(2.60\), but have \(1.95\)
and \(3.28\) total trades, respectively \cite{P}. Ada raised this pattern and Emmy checked the
figures. The trade counts compare two contract formats; they are not a comparison with no contract.
Q also reports that agents on the opposite side of the market can anchor on an opening offer
\cite{Q}. These observations motivate the proposed test, but do not show that plans cause fewer
trades in Q.

\section{Objections and unresolved points}

The peers objected that a bilateral contract could change Q's matching rules and make an apparent
improvement hard to interpret. Their proposed test keeps all execution in the existing book. They
also separated realised surplus, trade count, and price dispersion: more trades need not improve
price discovery, and dispersion based on very few trades may be misleading. Since all agents share
one book, Ada proposed assigning display formats to whole market runs and pairing runs with the same
reservation schedule and agent selection seed.

Emmy found a further comparison problem in P. Its natural language condition retains the negotiation
transcript, while its structured condition shows a compiled contract. The conditions differ in
context length and in how terms are interpreted, as well as in presentation. To isolate
presentation, the peers proposed showing identical verified terms in concise prose or structured
fields at the same point in each agent's context, with the full transcript tested separately.

Ada suggested checking whether agents wait for a planned partner after that partner has traded. Emmy
then corrected an information error: Q's agent-facing history does not name the traders in accepted
offers. An analyst could identify a partner's exit from the simulation state, but an agent would not
ordinarily know it. Continued reliance after an explicitly stated plan expiry is a cleaner
diagnostic. Even then, refusing a profitable offer can be reasonable if a better offer is expected;
the comparison must be between assigned conditions, not a verdict on one decision.

The ledger's prior-work check also limits any novelty claim. It identifies Hantel's study of a
seller anchor in an LLM market \cite{Hantel2026} and Agrawal and colleagues' study of communication
and collusion in double auctions \cite{Agrawal2025}. The peers did not independently assess those
works. Whether the narrower comparison of two displays of the same self-negotiated plan is new
remains unverified.

\section{Why the thread stopped}

The verifier paused the thread after one round. The supported conclusion is the boundary result: Q
already enforces crossing trades, and a payment tied to one trade is equivalent to changing its
price. The proposed effect of plan presentation remains a hypothesis. P changes more than
presentation across its contract conditions, and neither paper tests the effect in Q's auction.

The smallest step that would restart the thread is the controlled Q-style comparison the peers
specified. Form plans under one fixed protocol, then assign whole market runs to see the same terms
as concise prose or compact structured fields. Keep Q's order book and executable rights fixed, and
compare with an unmodified baseline. Measure realised surplus, crossing behaviour, missed gains at
the round limit, trade count, and price dispersion separately. Until such a test exists, the
direction and size of any effect remain unknown.

\bibliographystyle{plainurl}
\bibliography{references}

\end{document}
